The validated section JSON objects produced by the constrained leaf-generation stage are materialized into a final LaTeX manuscript and an accompanying BibTeX file through a deterministic assembly pipeline. The assembly pipeline performs three core functions: (1) structural mapping, (2) bibliography synthesis, and (3) provenance export and audit support.

Structural mapping translates JSON fields into canonical LaTeX constructs. For each section object the assembler emits a labeled LaTeX fragment using the section-level metadata (e.g., \texttt{title} $\rightarrow$ \verb|\section| / \verb|\subsection|), injects the validated plain-text body into a controlled environment, and attaches any supplied LaTeX snippets from the \texttt{latex\_snippets} field. Cross-reference metadata (the node identifier and dependency list) is rendered as stable \verb|\label| commands and used to populate an internal table of contents so the compiled document reflects the document-graph ordering. Transformations are recorded so that the mapping is reversible: the assembler preserves the originating section node identifier in an auxiliary LaTeX comment for each fragment, enabling later correspondence between PDF content and the source document graph.

Bibliography synthesis converts structured citation entries and per-claim source pointers contained in the section JSON (for example, \texttt{sources} and \texttt{evidence\_pointers}) into BibTeX records in \texttt{refs.bib} and into LaTeX citation tokens embedded at claim sites. The assembler uses structured citation metadata (author, title, venue, year, identifier when available) to generate a consistent BibTeX key and to canonicalize identifiers so as to avoid duplicate entries. Where a section supplies only an opaque source pointer, that pointer is preserved verbatim both as a BibTeX comment and as an inline audit footnote in the compiled document; the footnote text is generated from the \texttt{evidence\_pointers} field so that every cited claim in the human-readable PDF is associated with the original pointer returned by the retrieval stage.

Provenance export and audit support produce machine-parseable artifacts that enable automated verification and human review. For each assembled manuscript the pipeline writes a JSON manifest (for example, \texttt{manifest.json}) that lists, per section and per atomic claim, the originating node id, the exact byte offsets of the claim within the section body, the set of attached evidence pointers, and a checksum of the section JSON used for generation. This manifest, together with the original section JSON objects and \texttt{refs.bib}, forms an audit bundle that tools or reviewers can use to trace any assertion in the PDF back to its source pointers and to re-run validation steps deterministically; prior work motivates the need for black-box, machine-parseable provenance artifacts to support reliable auditing of generated content \cite{web_zhu_2025_auditing_data_provenance_in_real_world_text_to_i}.

Validation is applied at two stages of assembly. Prior to conversion, the assembler requires successful JSON-schema validation of each section and rejects or flags sections that fail required-field checks (missing \texttt{title}, missing \texttt{body}, or malformed \texttt{sources}). After conversion, a LaTeX sanity pass verifies that all generated \verb|\cite| and \verb|\label| tokens are resolved and that every claim annotated with evidence pointers has at least one corresponding entry in the manifest. Any unresolved citations, duplicate BibTeX keys, or missing provenance entries are surfaced in an assembly log so they can be inspected and corrected manually.

Practical considerations and common failure modes are recorded alongside each assembly run. Typical examples include duplicate or malformed source pointers that prevent canonical BibTeX generation, extracted metadata with insufficient fields for canonicalization, and sections that deliberately omit claims when supporting evidence is unavailable; such cases are preserved in the manifest as explicit "missing evidence" annotations to maintain an auditable history. Finally, we note a limitation that affects the practical completeness of the audit trail: its usefulness is bounded by retrieval coverage and by the fidelity of extracted metadata from source documents (general background knowledge).